\documentclass{resume}

\usepackage[left=0.75in,top=0.1in,right=0.75in,bottom=0.6in]{geometry}
\usepackage[utf8]{inputenc}
\usepackage{hyperref}
\usepackage{graphicx}
\usepackage{float}
\usepackage[absolute,overlay]{textpos}

%\setlength{\TPHorizModule}{1mm}
%\setlength{\TPVertModule}{1mm}
\graphicspath{ {./images/} }
\inputencoding{utf8}

\def\nameskip{\bigskip}
\def\sectionskip{\medskip}

\begin{document}

   %\begin{textblock}{20}(171,10)
   %   \includegraphics{java-certified-programmer}
   %\end{textblock}

   %\begin{textblock}{20}(171,30)
   %   \includegraphics{java-certified-web}
   %\end{textblock}

  \begin{center}
    {\fontsize{16}{18} \bf Murilo Cadanus da Costa} \\\\
  \end{center}
  \begin{center}
    {\fontsize{12}{14} 3803, Rue Nilo Peçanha} \\ 
    {\fontsize{12}{14} Curitiba, Paraná - 82120-440 - Brésil} \\
    {\fontsize{12}{14} murilocadanus@gmail.com} \\
    {\fontsize{12}{14} skype: murilocadanus}
  \end{center}

  \begin{rSection}{Sommaire}
  \end{rSection}

  \begin{itemize}
     \item Ingénieur de Logiciels détenant 12+ années d’expérience.
     \item Excellente connaissance de la programmation orientée objets avec C/C++/Objective-C et Java.
     \item Riche expérience dans le domaine de développement des logiciels multiplateforme%: Desktop, Mobile et Web.
     \item Capable d'apprendre et de s'adapter facilement.
     \item Français (avancé), Anglais (avancé) et Portugais (courant).
  \end{itemize}

  \begin{rSection}{Formation}
    {\bf Spécialisation en Informatique - Programmeur de Jeux Vidéo } \hfill {\em 2012} \\
    {PUC-PR, Curitiba - PR - Brésil} \\ 

    {\bf Baccalauréat en Informatique } \hfill {\em 2008} \\
    {Université Positivo, Curitiba - PR - Brésil}   
  \end{rSection}
  
  \begin{rSection}{Technologies Maitrisées}
    \begin{rSubsection}{}{}{}{}
      \item {\bf Math:} Systèmes de coordonnées cartésiennes, vecteurs, matrices.
      \item {\bf Concepts:} Design Patterns, OOP, UML, Scrum, Agile, SOA, BPM, TDD, BDD.
      \item {\bf Langages de programmation:} C/C++/Objective-C, Java/C#, Lua/Ruby.
      \item {\bf Frameworks/APIs - Graphique:} OpenGL, Direct3D, GLSL, HLSL.
      \item {\bf Frameworks/APIs/Utils - C/C++:} C standard library, Sockets, p\_threads, CMake, QMake, GDB, Valgrind.
      \item {\bf Frameworks/APIs/Utils - Java EE:} Servlets, Spring, EJB, MDB, JMS, JPA, JDBC, Web Services (SOAP et REST).
      \item {\bf Frameworks/APIs/Utils - Objective-C:} UIKit, Core Graphics, Core Animation, CoreData, CocoaPods.
      \item {\bf Engines:} Unity, UnrealEngine, Seed Framework.
      \item {\bf Bases de données}: MySQL, PostgreSQL, Oracle.
      \item {\bf Source Control Management Systems}: SVN, Git.
      \item {\bf Integrated Development Environments}: Qt Creator, XCode, Eclipse et Visual Studio.
      \item {\bf Modeling et Design}: Enterprise Architect, IBM Rational Software Architect.
    \end{rSubsection}
  \end{rSection}
  
  \begin{rSection}{Expérience Profissionnelle}

    % Sascar
    \begin{rSubsection}{\bf Architecte et programmeur}{Jan. 2015 - Actuelle}{Sascar}{}
      La compagnie Sascar est une entrepise de ingénierie qui devellope d'applications de tracking. Elle est à Curitiba au Brésil. J'ai travaillé comme architecte de solutions et programmeur, mon travail était la conception, le développement et la recherche de meilleures solutions pour communication avec les equipment de localisation.
    \end{rSubsection}
    \begin{rSubsection}{Responsabilités:}{}{}{}
      \item Crée l’architecture des applications de traitement de commandes d'equipment;
      \item Analysé, programmé et testé le code en C/C++;
      \item Évalué et testé nouvelles tecnologies;
    \end{rSubsection}
    \begin{rSubsection}{Projets significatifs:}{}{}{}
      \item Projet LibI9 - La lib est la engine responsable comme base pour generalisation des components: Log, Threads, Sockets, IO. J’ai développé la lib pour GNU/Linux en \textbf{C/C++} avec \textbf{QTCreator}. \\
      \item Projet TPMS - Le TPMS est un traqueur de pneu utilise pour savoir la pression et temperature. J’ai développé la application de communication utilisant \textbf{sockets} \textbf{TCP/IP}, \textbf{p\_threads} en \textbf{C/C++} avec \textbf{QTCreator}. \\
      \item Projet Bluetec - Le Bluetec est une application de tracking de véhicules . J’ai développé le parse de commands en \textbf{C/C++} avec \textbf{QTCreator}. \\
    \end{rSubsection}

     {\fontsize{8}{9}\selectfont \textbf{Environnement technique:} C/C++, Sockets, p\_threads, QMake, Git, Git flow.}

    % PlugApps
    %\begin{rSubsection}{\bf Architecte et programmeur}{Jan. 2015 - Actuelle}{PlugApps}{}
    %  La compagnie PlugApps est une entrepise startup qui devellope d'applications mobile. Elle est à Curitiba au Brésil. J'ai travaillé comme architecte de solutions mobile et programmeur, mon travail était la conception, le développement et la recherche de meilleures solutions pour les applications iOS.
    %\end{rSubsection}
    %\begin{rSubsection}{Responsabilités:}{}{}{}
    %  \item Crée l’architecture des applications iOS;
    %  \item Analysé, programmé et testé le code en Objective-C;
    %  \item Évalué et testé nouvelles tecnologies;
    %\end{rSubsection}
    %\begin{rSubsection}{Projets significatifs:}{}{}{}
    %  \item Projet WHATSEE (\url{http://whatsee.photo}) - Le WhatSee est une reseau social de photo. J’ai développé l'application mobile pour iOS en \textbf{Objective-C} avec \textbf{Xcode}. \\
    %  \item Projet Reverse (\url{http://reverse.best}) - Le Reverse est une service d'achat où les annonce sont pour acheter de produit et non pour vendre. J’ai développé le application mobile pour iOS en \textbf{Objective-C} avec \textbf{Xcode}. \\
    %\end{rSubsection}

     %{\fontsize{8}{9}\selectfont \textbf{Environnement technique:} Xcode, Objective-C, Cocoapods, UIKit, Push notification, Core data, AsyncTasks(GCD et Operation Queue), Reactive Cocoa, AFNetworking, SDWebImage, REST Web services, Postman, XMPPFramework, Expecta, OCMock, Specta, Git, Git flow.}

    % Winterlabs
    \begin{rSubsection}{\bf Architecte et programmeur}{Déc. 2010 - Déc. 2014}{Winterlabs Technologies}{}
      La Winterlabs est une studio de development des jeux videos mobile. J'ai travaillé comme architecte et programmeur, mon travail était la conception, le développement et la recherche de meilleures solutions pour les jeux videos d'entreprise.
    \end{rSubsection}
    \begin{rSubsection}{Responsabilités:}{}{}{}
      \item Crée l’architecture des jeux videos;
      \item Analysé, programmé et testé le code en C/C++/Objective-C et Java;
      \item Effectué le déploiement, la gestion et le support des jeux videos;
      \item Évalué et testé nouvelles tecnologies;
    \end{rSubsection}
    \begin{rSubsection}{Projets significatifs:}{}{}{}
      \item Projet SKY FIGHTERS (\url{http://itunes.apple.com/us/app/skyfighters-hd/id419328364}) - Sky Fighters est un jeu de carte et stratégie pour \textbf{iOS}. Je l’ai développé en \textbf{Objective-C} avec \textbf{Xcode}, framework \textbf{cocos2d}. \\
      \item Projet RESTA1 (\url{http://itunes.apple.com/br/app/resta1/id322697159}) – Jeu de plateau classique pour \textbf{iOS} dans lequel le joueur doit éliminer le nombre maximum de pièces du plateau. Je l’ai développé en \textbf{Objective-C} avec \textbf{Xcode}, \textbf{cocos2d}.
    %\item Projet ETERNAL LEGENDS (\url{http://youtu.be/LW1yF8D-CKU}) - Le Eternal Legends est un jeu vidéo de carte en 3D pour le Mac OS X et le Windows. Le joueur doit utiliser ses cartes pour détruire l'adversaire. Je l’ai développé en \textbf{C++} avec \textbf{Xcode} et \textbf{VisualStudio}, l’architecture utilise \textbf{SDL} pour contrôler les actions d’entrée et le son. Pour le rendu j’ai utilisé \textbf{OpenGL}. Les agents de IA ont été développés et exposés comme scripts \textbf{Lua} pour faciliter la modularisation. J’ai utilisé \textbf{Git} pour la gestion du code source. \\
    %\item Projet BERMUDA TRIANGLE (\url{https://bermudatriangle.herokuapp.com}) - Le Bermuda Triangle est un jeu 2D procédural, où le joueur doit collecter pièces pour ponctuer. Je l'ai développé en \textbf{C\#} avec \textbf{VisualStudio}, \textbf{Unity3D} et \textbf{NGUI} plugin. J'ai integré avec le Apple IAP pour vendre les items en jeu et Apple Game Centre pour le highscores. J'ai utilisé le \textbf{Git} pour la gestion du code source. \\
    %\item Projet CORE RUNNER (\url{https://itunes.apple.com/us/app/core-runner/id821017105}) - Le Core Runner est un jeu 3D procédural, où le joueur doit courir et collecter cubes pour ponctuer. Je l'ai développé en \textbf{C\#} avec \textbf{MonoDevelop}, \textbf{Unity3D} et \textbf{NGUI} plugin. J'ai integré avec le Apple IAP pour vendre les items en jeu et Apple Game Centre pour le highscores. J'ai utilisé le \textbf{Git} pour la gestion du code source. \\
    %\item Projet STAP – C'est un simulateur de tir virtuel pour les armes portables. J'ai développé le module qui contrôle la piste de tir utilisant \textbf{C\#} avec \textbf{VisualStudio}, J'ai utilisé l'\textbf{Unity3D} pour rendre les scènes du stand de tir et \textbf{NGUI} plugin pour l'UI, \textbf{Crystal Reports} pour développer les rapports sur les sessions de tir et \textbf{MySQL} comme sauvegarde.  J'ai utilisé le \textbf{Mercurial} pour la gestion de versions du source code. \\
    %\item Projet SALVA CARGAS (\url{http://salvacargas.herokuapp.com}) - Le Salva Cargas est un jeu éducatif avec l'objective de développer le raisonnement mathématique. Le jouer doit utiliser un canon et deux bateaux pour récupérer des charges qui volent sur une île. Je l'ai développé en \textbf{C\#} avec \textbf{MonoDevelop}, \textbf{Unity3D} et \textbf{NGUI} plugin. J'ai utilisé le \textbf{Git} pour la gestion du code source. Le jeu a été compilé pour fonctionner avec la technologie Flash. \\
    %\item Projet MR FALL (\url{http://itunes.apple.com/br/app/mrfall/id420213660}) – Mr Fall est un jeu de plateforme en 2D pour \textbf{iOS}. Le jeu a 12 niveaux répartis dans trois différents paysages: montagne, forêt et cave. L'objectif est de recueillir tous les diamants dans le niveau et trouver la sortie. Je l’ai développé en \textbf{C\#} avec \textbf{MonoDevelop}, \textbf{Unity3D}, j'ai utilisé comme moteur physique le \textbf{PhysX}, j'ai intégré avec Mobclix, AdMob et Millenial Media network et \textbf{Git}. \\
    \end{rSubsection}

    {\fontsize{8}{9}\selectfont \textbf{Environnement technique:} Xcode, C, C++, Objective-C, Java, JEE 1.5, PostgreSQL, Cocos2D, SFML, SDL, OpenGL, GLSL, Git, Git flow.}

    %\begin{rSubsection}{Global GameJam, 7DRL}{2012, 2013 et 2014}{Programmeur}{Curitiba, PR}
    %\item Projet SEED FRAMEWORK(\url{https://github.com/fungos/seed/commits?author=murilocadanus}) - C'est un moteur de développement de jeux vidéo portable de source code ouvert. J'ai développé l'infrestucture multijoueur avec UDP Sockets, l'infrastucture de Shaders, l'infrastructure de pathfind A*. utilisant \textbf{C++} avec l'IDE \textbf{QTCreator} et le \textbf{Git} pour la gestion de versions du source code. \\
    %\item Projet REAPERS (\url{https://github.com/ggj/reapers/commits?author=murilocadanus}) – C'est un jeu de plateforme 2D. J'ai développé utilisant \textbf{C++} avec l'IDE \textbf{QTCreator}, le moteur \textbf{Seed} et le \textbf{libRocket} API pour l'UI.  J'ai utilisé le \textbf{Git} pour la gestion de versions du source code. \\
    %\item Projet OPTIMALUS (\url{https://github.com/ggj/optimalus/commits?author=murilocadanus}) – C'est un jeu de style Rogue Like. J'ai utilisé \textbf{C++} avec l'IDE \textbf{QTCreator}, le moteur \textbf{Seed} et le \textbf{libRocket} API pour l'UI.  J'ai utilisé le \textbf{Git} pour la gestion de versions du source code. \\
    %\item Projet QUAKERL (\url{https://github.com/7drl/QuakeRL/commits?author=murilocadanus})– J'ai développé le module de pathfind utilisant \textbf{C++} avec l'IDE \textbf{QTCreator}, le moteur \textbf{Seed} et le \textbf{libRocket} API pour l'UI.  J'ai utilisé le \textbf{Git} pour la gestion de versions du source code. \\
    %\end{rSubsection}
    % GVT
    \begin{rSubsection}{\bf Analyste et programmeur}{Août 2008 - Déc. 2010}{GVT Telecom}{}
    La compagnie GVT(Global Village Telecom) est une grand entrepise de télécommunications brésilienne qui offre des services sur le téléphone fixe, le haut débit à la fois consommateurs et des entreprises, la télévision payante et la voix sur IP. GVT a été sur le marché depuis la fin de l'année 2000. J'ai travaillé comme analyste et programmeur, mon travail était la conception, le développement et la recherche de meilleures solutions pour les applications d'entreprise.
    \end{rSubsection}
    \begin{rSubsection}{Responsabilités:}{}{}{}
      \item Analysé, programmé et testé le code en Java;
      \item Développé une nouvelle interface de communication avec le nouvel équipement ADSL
      \item Effectué le déploiement, la gestion et le support de la solution;
    \end{rSubsection}
    \begin{rSubsection}{Projets significatifs:}{}{}{}
    \item Projet "Nova Arquitetura OSS” – Modification de nouveaux équipements pour augmenter la vitesse du produit ADSL. J'ai développé une nouvelle interface de communication avec le nouvel équipement ADSL, j'ai utilisé \textbf{JEE 1.5} avec \textbf{Oracle 10g} et le serveur \textbf{BEA Weblogic application server}, \textbf{BEA Aqualogic Service Bus}, \textbf{EJB3}, \textbf{JPA}, \textbf{JAX-WS 2.0}, \textbf{JAXB} et \textbf{Eclipse}.\\
    \item Projet PD “Modificação Portal Integra” – Modification pour fournir de nouvelles informations à la clientèle. J'ai développé l'interface  utilisant des \textbf{portlets} avec \textbf{BPM} et l'architecture \textbf{SOA} utilisant \textbf{Savvion}, j'ai utilisé le serveur d'applications \textbf{BEA Weblogic Portal}, \textbf{BEA Aqualogic Service Bus}.\\
    \item Projet “PJ 759 – Nova Regulamentação de Call Centers”. Projet pour répondre à la nouvelle réglementation des centres d'appel. J‘ai spécifié le IT Solution avec \textbf{Enterprise Architec} et j'ai développé les \textbf{Stored Procedures} et \textbf{Triggers} en \textbf{PL/SQL}, \textbf{Oracle} comme \textbf{SGBDR} et \textbf{JSP}, \textbf{JSTL}, \textbf{Struts}, \textbf{EJB 2.1} sous \textbf{Oracle IAS} serveur d'application.\\
    \item Projet “PJ 679 - TT-Retail”. Projet avec l’objectif de manoeuvre des stations de base existantes. J’ai utilisé \textbf{BEA WebLogic Workshop 8}, \textbf{BEA AquaLogic Service Bus}, \textbf{WebLogic Integration (WLI)}, \textbf{JMS}, \textbf{Oracle AQ}, \textbf{MDB} et \textbf{BEA Weblogic Application Server 8}.\\
    \item Projet PJ 771 - WiseTools. Web services qui exposent les données contenues de la base du TBS/OSS. J’ai utilisé \textbf{JEE 1.5} avec \textbf{Oracle 10g} et \textbf{BEA Weblogic serveur d'application}, \textbf{BEA AquaLogic Service Bus}, \textbf{EJB3}, \textbf{JPA}, \textbf{JAX-WS 2.0}, \textbf{JAXB}, \textbf{Eclipse}, \textbf{Mylyn}. J'ai utilisé une méthode de developpement agile (\textbf{SCRUM}) avec l'outil \textbf{JIRA} et \textbf{Green Hopper} (agile plugin pour burndown chart).\\
    \item Projet “PJ 702 – Serviços de Gerência”. C’est un système SLA pour gérer les données et la qualité du transfert de voix des clients directs. J’ai utilisé les technologies \textbf{Servlet}, \textbf{JSP}, \textbf{JSTL}, \textbf{AJAX}, \textbf{EJB-doclet}, \textbf{EJB 2.1}, \textbf{Spring}, \textbf{JMS}, \textbf{MDB}, \textbf{Web Services}, \textbf{JasperReports}, \textbf{JFreeChart}, \textbf{ANT} et les design patterns, \textbf{Session Facade}, \textbf{DTO}, \textbf{Business Object}, \textbf{Service Locator}. J'ai utilisé le \textbf{BEA WebLogic Workshop 9.2} comme IDE, le \textbf{BEA Weblogic 9.2} comme serveur d'applications et \textbf{BEA AquaLogic Service Bus} comme \textbf{ESB}.\\
    \item Projet “Arquitetura Integração Contínua OSS” – Definition de l’architecture de développement d’intégration continue pour les projets de l'équipe OSS-GVT. J'ai utilisé les technologies: \textbf{SVN} pour contrôle de code source, \textbf{Maven} pour compiler le code source, le programme \textbf{NEXUS} pour contrôle d'artefacts, \textbf{BAMBOO} pour gérer et planifier les actions de code source, \textbf{JIRA} comme Issue Tracker, \textbf{JUnit} pour tester le code, \textbf{Selenium} pour faire des tests unitaires automatiques et les plugins \textbf{Checkstyle}, \textbf{PMD}, \textbf{STATSVN}, \textbf{Cobertura}, \textbf{FindBugs} pour faire des rapports sur la qualité du code.
    \end{rSubsection}
    
    {\fontsize{8}{9}\selectfont \textbf{Environnement technique:} Sparx Enterprise Architect, Eclipse, Java, JEE 1.5, Servlets, JSP, JSF, EJB (Session Beans et Entity Beans), JPA, MDB, JMS, Web Services, SOA, BPM, Savvion, ESB, Oracle Weblogic, AquaLogic Service Bus, Oracle AQ, Oracle 10g, PostgreSQL, Ant, Maven, Atlassian JIRA, SVN.}

    \begin{rSubsection}{\bf Analyste et programmeur}{Janv. 2006 - Mai. 2008}{HSBC Banque}{}
      HSBC est un groupe bancaire international britannique présent au Brésil depuis 1997. Le department de informatique est à Curitiba au Brésil. J'ai travaillé comme analyste et programmeur, mon travail était la conception, le développement et la recherche de meilleures solutions pour les applications d'entreprise.
    \end{rSubsection}
    \begin{rSubsection}{Responsabilités:}{}{}{}
      \item Analysé, programmé et testé le code en Java;

    \end{rSubsection}
    \begin{rSubsection}{Projets significatifs:}{}{}{}
    \item Projet “HOB-PWS” - Système développé pour les clients sans compte bancaire; modules Intranet et Internet. J’ai utilisé (Intranet): pattern \textbf{MVC} (\textbf{ITA Framework} / \textbf{MVC2}), intégration avec \textbf{CICS} \textbf{IBM MQSeries} les technologies (Internet): \textbf{Servlets}, \textbf{XML} et \textbf{XSL}. \textbf{WSAD}, \textbf{Websphere 5.1}.\\
    \item Projet “CDCI”. Système de simulation de financement direct pour le client. J’ai fait la spécification de l’application utilisant le \textbf{IBM RSA}. J’ai développé dans le modèle \textbf{MVC} avec \textbf{Eclipse-RCP}, \textbf{SWT}, la sérialisation d'objets comme couche de sauvegarde.\\
    \item Projet “HOB-PG-TIT”. Système de paiement de titres utilisés par les clients qui accèdent au site internet de la banque HSBC. Je l’ai développé avec \textbf{Servlets}, \textbf{XML} et \textbf{XSL}, Architecture \textbf{SOA} avec le \textbf{BPM} utilisant \textbf{RTP} (processeur de transactions à distance), l'intégration avec l'héritage (\textbf{CICS}), \textbf{IBM MQSeries}, \textbf{WSAD} (WebSphere Application Developer), \textbf{WebSphere 5.1}.\\
    \item Projet “GRH - Système de ressources humaines. Je l’ai développé avec le modèle \textbf{MVC} (\textbf{Framework ITA} / \textbf{MVC2}) avec \textbf{JDBC} et \textbf{Sybase}. \textbf{WSAD}; \textbf{Websphere 5.1}
    \end{rSubsection}

    {\fontsize{8}{9}\selectfont \textbf{Environnement technique:} Java, JEE 1.5, Servlets, JSP, XML, XSL, JDBC, Sybase, IBM Websphere, IBM WSAD, IBM RSA, IBM RSM, Eclipse, SOA, BPM, CVS, SVN.}
  \end{rSection}

  \begin{rSection}{Certifications}
    {\bf Sun Microsystems, Inc.} \hfill {\em Mai 2005} \\ 
    {Sun Certified Programmer for the Java 2 Platform 1.4} \\

    {\bf Sun Microsystems, Inc.} \hfill {\em Mai 2008} \\ 
    {Sun Certified Web Component Developer for the Java 2 Platform Enterprise Edition 1.5} \\
  \end{rSection}
  
\end{document}
